\chapter{Elicitación de requisitos}

\begin{quotation} [Empresario] {Ray Kroc (1902--1984)}
Los dos requisitos más importantes para el éxito principal son los siguientes: primero, estar en el lugar correcto en el momento adecuado, y en segundo lugar, hacer algo al respecto.
\end{quotation}

\begin{abstract}
Para que los proyectos lleguen a buen puerto se han de tener claro desde un principio qué se consideran requisitos indispensables para que el proyecto sea un éxito.
\end{abstract}

\section{Planificación de la gestión de los requisitos}

La elicitación de requisitos es uno de los mayores retos a la hora de afrontar un proyecto, para esto es necesario tener en conocimiento el dominio del proyecto así como el problema que se trata de solucionar con este, así como a sus usuarios potenciales. Esto nos ayudará a tener claro las necesidades y expectativas tanto de los usuarios como del proyecto en sí.

Previa a la propia elicitación de los requisitos es necesario aclarar que dispondremos de cinco tipo de requisitos, los requisitos funcionales, los requisitos no funcionales, las reglas de negocio, los requisitos de información y los requisitos de calidad.

\subsection{Definición de los requisitos}
Todos los requisitos seguirán el siguiente formato: \emph{RX-001}, en el que la X se corresponderá con la inicial o iniciales del tipo de requisito que se especifique, el serial numérico que le seguirá será un identificador numérico único para ese tipo de requisito. A este serial le acompañará una breve descripción del requisito

\subsection{Métricas para controlar el cumplimiento de los requisitos}

\begin{itemize}

\item Tiempo: Una de las principales métricas de un proyecto, especialmente cuando tenemos un hito crítico en una fecha concreta o una fecha fin inamovible. La gestión del tiempo se deberá hacer de manera frecuente durante todo el proyecto y esto nos va a permitir realizar predicciones a futuro. Los principales indicadores de tiempo son el \textbf{SV} (Schedule Variance) y el \textbf{SPI} (Schedule Performance Index).

\item Satisfacción del cliente: Se refiere al grado en que los resultados del proyecto cumplen o superan las expectativas. Esto abarca la calidad de las entregas, la experiencia general del cliente, el servicio al cliente, y la comunicación entre las partes interesadas internas y externas a lo largo del ciclo de vida del proyecto.

\item Presupuesto: El presupuesto es lo que costará desarrollar el proyecto. El presupuesto está intrínsecamente ligado a los otros criterios mencionados anteriormente. Si los costos reales comienzan a superar los gastos planificados, se deberá ajustar el tiempo, el alcance y/o la satisfacción del cliente para alcanzar el punto de equilibrio.
Para ellos emplearemos \textbf{KPIs} (indicadores claves de rendimiento) específicos para rastrear que pueden ayudarle a mantener su proyecto dentro del presupuesto incluyen:

\begin{itemize}

\item Valor planificado
\item Costo real
\item Valor ganado
\item Índice de rentabilidad: Valor ganado (EV) / coste real (AC)

\end{itemize}

\end{itemize}

\subsection{Análisis del impacto y niveles de autorización de un cambio}

Se deberá realizar una reunión con la totalidad de los desarrolladores, para aprobar, dejar pendiente o rechazar los cambios solicitados siguiendo el procedimiento definido anteriormente. Estas decisiones no podrán ser tomadas directamente por el Proyect Manager puesto que pueden afectar las prioridades del proyecto así como el alcance de este, resultando en retrasos. Por ello la unanimidad de los desarrolladores será necesaria para la aprobación de un cambio en los requisitos, puesto que estos tienen el control total del proyecto.


\subsection{Priorización de los requisitos del proyecto}
Para la elección de requisitos, debemos establecer unas prioridades sobre el alcance del proyecto, es decir, que es prioritario para el proyecto, aquellos objetivos imprescindibles para el cumplimiento de las necesidades del cliente, y que otros requisitos añaden valor al proyecto para algunos de los interesados, pero no son imprescindibles para el negocio. Una de las técnicas más empleadas en la priorización de requisitos es la conocida como técnica \textbf{MoSCoW \cite{moscow}}. Esta técnica nos ayudará a separar los requisitos en cuatro grandes categorías:

\begin{itemize}

\item \textbf{Must (M)}: Requisitos fundamentales y obligatorios para satisfacer las necesidades del negocio. Si estos no se cumplen, se verá 
\item \textbf{Should (S)}: Requisitos que deberían ser cumplidos en la medida de lo posible. El éxito del proyecto no dependerá de estos, aunque podría reducir su alcance. Estos requisitos pueden ser satisfechos mediante soluciones temporales, o en coso de que existiera una causa justificada, se podría prescindir de estos.
\item \textbf{Could (C)}: Requisitos que son interesantes que cumpla el servicio. Se trata de requisitos adicionales que se implementarán en el caso de disponer de tiempo y presupuesto para ello. Estos requisitos mejorarían el rendimiento del servicio pero podrían ser eliminados fácilmente.
\item \textbf{Won´t (W)}: Requisitos que se ha decidido no implementar de momento, pero que serán tomados en cuenta en el futuro con el objetivo de mejorar el servicio o producto. 

\end{itemize}

\section{Documentación de requisitos}
A continuación, definimos los requisitos determinantes para la correcta evolución del proyecto.

\subsection{Reglas de negocio}

Las reglas de negocio son el tipo de requisito más importante de todos puesto que estarán determinados por el tipo de proyecto a desarrollar, entonces en el contexto del proyecto de T-Planifica, las reglas de negocio serán:

\begin{table*}[htb]
	\centering
	\begin{coolTable}{p{4cm}p{\textwidth-6cm}}{2}{RN-001}
\textbf{DEFINICIÓN}& \textbf {Como} estudiante, \textbf {quiero} tener que realizar el menor número de permutas posibles con otros estudiantes,  \textbf {para} entregar el menor número de documentos posibles en la secretaría.\\
\midrule
\textbf{PRIORIDAD}&M\\
\midrule
\textbf{CRITERIOS \newline DE ACEPTACIÓN}& Minimizar las permutas entre alumnos maximizando los cambios de grupo. \\
	\end{coolTable}
	\caption{Regla de negocio 001\label{tab:RN-001}}
\end{table*}

%-------------------------------------------------------------------------------------------------------
\begin{table*}[htb]
	\centering
	\begin{coolTable}{p{4cm}p{\textwidth-6cm}}{2}{RN-002}
\textbf{DEFINICIÓN}& \textbf {Como} estudiante, \textbf {quiero} que el documento de permuta se genere en el servidio,  \textbf {para} evitar hacerlo yo de manera incorrecta.\\
\midrule
\textbf{PRIORIDAD}&M\\
\midrule
\textbf{CRITERIOS \newline DE ACEPTACIÓN}& Cuando se acepte una permuta por ambas partes se podrá descargar un documento con los campos rellenos para la posterior firma por ambas partes y la entrega. \\
	\end{coolTable}
	\caption{Regla de negocio 002\label{tab:RN-002}}
\end{table*}

%-------------------------------------------------------------------------------------------------------
\begin{table*}[htb]
	\centering
	\begin{coolTable}{p{4cm}p{\textwidth-6cm}}{2}{RN-003}
\textbf{DEFINICIÓN}& \textbf {Como} estudiante, \textbf {quiero} que se me proponga la mejor opción de permuta, \textbf {para} evitar la búsqueda del cambio más eficiente.\\
\midrule
\textbf{PRIORIDAD}&M\\
\midrule
\textbf{CRITERIOS \newline DE ACEPTACIÓN}& Se realicen la asignación de permutas obteniendo los mejores intercambios. \\
	\end{coolTable}
	\caption{Regla de negocio 003\label{tab:RN-003}}
\end{table*}

%-------------------------------------------------------------------------------------------------------

\begin{table*}[htb]
	\centering
	\begin{coolTable}{p{4cm}p{\textwidth-6cm}}{2}{RN-004}
\textbf{DEFINICIÓN}& \textbf {Como} estudiante, \textbf {quiero} conocer la trazabilidad del proceso de permutas,  \textbf {para} asegurar la transparencia en el etapas del proceso del intercambio de grupos.\\
\midrule
\textbf{PRIORIDAD}&M\\
\midrule
\textbf{CRITERIOS \newline DE ACEPTACIÓN}& Mostrar el estado en el que se encuentra la solicitud. \\
	\end{coolTable}
	\caption{Regla de negocio 004\label{tab:RN-004}}
\end{table*}

%-------------------------------------------------------------------------------------------------------

\begin{table*}[htb]
	\centering
	\begin{coolTable}{p{4cm}p{\textwidth-6cm}}{2}{RN-005}
\textbf{DEFINICIÓN}& \textbf {Como} estudiante, \textbf {quiero} poder contactar con compañeros, \textbf {para} ponerme en contacto en el proceso de formalización de la permuta.\\
\midrule
\textbf{PRIORIDAD}&M\\
\midrule
\textbf{CRITERIOS \newline DE ACEPTACIÓN}& Mostrar los mensajes de los usuarios que están interacturando en la aplicación. \\
	\end{coolTable}
	\caption{Regla de negocio 005\label{tab:RN-005}}
\end{table*}

%-------------------------------------------------------------------------------------------------------
\begin{table*}[htb]
	\centering
	\begin{coolTable}{p{4cm}p{\textwidth-6cm}}{2}{RN-006}
\textbf{DEFINICIÓN}& \textbf {Como} estudiante, \textbf {quiero} poder contactar con el administrador de la aplicación, \textbf {para} consultar dudas sobre el proceso de formalización de la permuta.\\
\midrule
\textbf{PRIORIDAD}&M\\
\midrule
\textbf{CRITERIOS \newline DE ACEPTACIÓN}& Solucionar las dudas sugeridas por los usuarios de la aplicación. \\
	\end{coolTable}
	\caption{Regla de negocio 006\label{tab:RN-006}}
\end{table*}

%-------------------------------------------------------------------------------------------------------

\begin{table*}[htb]
	\centering
	\begin{coolTable}{p{4cm}p{\textwidth-6cm}}{2}{RN-007}
\textbf{DEFINICIÓN}& \textbf {Como} estudiante, \textbf {quiero} poder generar un ticket desde la aplicación o el bot, \textbf {para} reportar problemas técnicos rápidamente.\\
\midrule
\textbf{PRIORIDAD}&M\\
\midrule
\textbf{CRITERIOS \newline DE ACEPTACIÓN}& El sistema debe generar un identificador único para cada ticket. Dicho ticket debe ser visible en el historial del usuario inmediatamente después de crearlo. \\
	\end{coolTable}
	\caption{Regla de negocio 007\label{tab:RN-007}}
\end{table*}

%-------------------------------------------------------------------------------------------------------

\begin{table*}[htb]
	\centering
	\begin{coolTable}{p{4cm}p{\textwidth-6cm}}{2}{RN-008}
\textbf{DEFINICIÓN}& \textbf {Como} estudiante, \textbf {quiero} recibir notificaciones sobre actualizaciones de mi ticket, \textbf {para} estar informado del estado de mi incidencia.\\
\midrule
\textbf{PRIORIDAD}&M\\
\midrule
\textbf{CRITERIOS \newline DE ACEPTACIÓN}& Las notificaciones deben enviarse cuando: \\
\begin{itemize}
    \item Se cambie el estado del ticket (por ejemplo, "En proceso" o "Resuelto").
    \item Se asigne un responsable al ticket.
\end{itemize}
	\end{coolTable}
	\caption{Regla de negocio 008\label{tab:RN-008}}
\end{table*}

%-------------------------------------------------------------------------------------------------------

\begin{table*}[htb]
	\centering
	\begin{coolTable}{p{4cm}p{\textwidth-6cm}}{2}{RN-009}
\textbf{DEFINICIÓN}& \textbf {Como} administrador, \textbf {quiero} registrar las asignaturas y asociar los diferentes grupos y proyectos docentes correspondientes, \textbf {para} que el estudiantado pueda elegir el grupo.\\
\midrule
\textbf{PRIORIDAD}&M\\
\midrule
\textbf{CRITERIOS \newline DE ACEPTACIÓN}& Mostrar las asignaturas y los grupos creados. \\
	\end{coolTable}
	\caption{Regla de negocio 009\label{tab:RN-009}}
\end{table*}

%-------------------------------------------------------------------------------------------------------
\begin{table*}[htb]
	\centering
	\begin{coolTable}{p{4cm}p{\textwidth-6cm}}{2}{RN-010}
\textbf{DEFINICIÓN}& \textbf {Como} administrador, \textbf {quiero} gestionar las preguntas frecuentes almacenadas del bot, \textbf {para} actualizar el contenido según las necesidades de los usuarios.\\
\midrule
\textbf{PRIORIDAD}&M\\
\midrule
\textbf{CRITERIOS \newline DE ACEPTACIÓN}& \\
\begin{itemize}
    \item Las preguntas frecuentes deben estar organizadas por categorías para facilitar la navegación.
    \item El administrador debe poder añadir, editar o eliminar preguntas frecuentes desde un panel de control.
\end{itemize}
	\end{coolTable}
	\caption{Regla de negocio 010\label{tab:RN-010}}
\end{table*}

%-------------------------------------------------------------------------------------------------------

\begin{table*}[htb]
	\centering
	\begin{coolTable}{p{4cm}p{\textwidth-6cm}}{2}{RN-011}
\textbf{DEFINICIÓN}& \textbf {Como} administrador, \textbf {quiero} poder categorizar y priorizar los tickets recibidos, \textbf {para} asegurar que los problemas críticos sean atendidos primero.\\
\midrule
\textbf{PRIORIDAD}&M\\
\midrule
\textbf{CRITERIOS \newline DE ACEPTACIÓN}& \\
\begin{itemize}
    \item Cada ticket debe tener un nivel de prioridad asignado: Bajo, Medio o Alto.
    \item Los responsables deben poder cambiar la prioridad de un ticket si la situación lo requiere.
\end{itemize}
	\end{coolTable}
	\caption{Regla de negocio 011\label{tab:RN-011}}
\end{table*}

%-------------------------------------------------------------------------------------------------------

\clearpage
\subsection{Requisitos de información}

\begin{table*}[htb]
	\centering
	\begin{coolTable}{p{4cm}p{\textwidth-6cm}}{2}{RI-001}
\textbf{DEFINICIÓN}& \textbf {Como} estudiante, \textbf {quiero} conocer mis datos en la aplicación,  \textbf {para} garantizar que los datos almacenados en la aplicación son correctos.\\
\midrule
\textbf{PRIORIDAD}&M\\
\midrule
\textbf{CRITERIOS \newline DE ACEPTACIÓN}& Se muestra en una pestaña los datos almacenados en la aplicación.\\
	\end{coolTable}
	\caption{Requisito de información 001\label{tab:RI-001}}
\end{table*}


\begin{table*}[htb]
	\centering
	\begin{coolTable}{p{4cm}p{\textwidth-6cm}}{2}{RI-002}
\textbf{DEFINICIÓN}& \textbf {Como} estudiante, \textbf {quiero} introducir la información adicional que sea necesaria, \textbf {para} realizar el proceso de permuta.\\
\midrule
\textbf{PRIORIDAD}&M\\
\midrule
\textbf{CRITERIOS \newline DE ACEPTACIÓN}& Se habilita un formulario para que el usuario pueda introducir la información requerida.\\
	\end{coolTable}
	\caption{Requisito de información 002\label{tab:RI-002}}
\end{table*}

\begin{table*}[htb]
	\centering
	\begin{coolTable}{p{4cm}p{\textwidth-6cm}}{2}{RI-003}
\textbf{DEFINICIÓN}& \textbf {Como} estudiante, \textbf {quiero} que se muestren los cursos de mi grado,  \textbf {para} filtrar entre los que estoy matriculado.\\
\midrule
\textbf{PRIORIDAD}&M\\
\midrule
\textbf{CRITERIOS \newline DE ACEPTACIÓN}& Se visualizan los cursos a los que se encuentra matriculado el estudiante.\\
	\end{coolTable}
	\caption{Requisito de información 003\label{tab:RI-003}}
\end{table*}

\begin{table*}[htb]
	\centering
	\begin{coolTable}{p{4cm}p{\textwidth-6cm}}{2}{RI-004}
\textbf{DEFINICIÓN}& \textbf {Como} estudiante, \textbf {quiero} que se muestren las asignaturas que pertenecen a un curso,  \textbf {para} filtrar las asignaturas que pertenecen a dicho curso.\\
\midrule
\textbf{PRIORIDAD}&M\\
\midrule
\textbf{CRITERIOS \newline DE ACEPTACIÓN}& Se listan todas las asignaturas que pertenecen al curso seleccionado por el estudiante.\\
	\end{coolTable}
	\caption{Requisito de información 004\label{tab:RI-004}}
\end{table*}

\begin{table*}[htb]
	\centering
	\begin{coolTable}{p{4cm}p{\textwidth-6cm}}{2}{RI-005}
\textbf{DEFINICIÓN}& \textbf{Como} estudiante, \textbf{quiero} poder consultar las permutas solicitadas por estudiantes, \textbf{para} buscar si existe algún intercambio que cumpla mis necesidades.\\
\midrule
\textbf{PRIORIDAD}&M\\
\midrule
\textbf{CRITERIOS \newline DE ACEPTACIÓN}& Se muestra un listado con las permutadas solicitadas por otros estudiantes.\\
	\end{coolTable}
	\caption{Requisito de información 005\label{tab:RI-005}}
\end{table*}

\begin{table*}[htb]
	\centering
	\begin{coolTable}{p{4cm}p{\textwidth-6cm}}{2}{RI-006}
\textbf{DEFINICIÓN}& \textbf{Como} estudiante, \textbf{quiero} poder marcar las permutas como favoritas, \textbf{para} encontrarlas de manera más ágil.\\
\midrule
\textbf{PRIORIDAD}&M\\
\midrule
\textbf{CRITERIOS \newline DE ACEPTACIÓN}& Se muestra un listado con las permutas marcadas como favoritas.\\
	\end{coolTable}
	\caption{Requisito de información 006\label{tab:RI-006}}
\end{table*}


\begin{table*}[htb]
	\centering
	\begin{coolTable}{p{4cm}p{\textwidth-6cm}}{2}{RI-007}
\textbf{DEFINICIÓN}& \textbf{Como} estudiante, \textbf{quiero} rellenar el formulario con los datos requeridos, {para} cumplimentar la emisión de la solicitud.\\
\midrule
\textbf{PRIORIDAD}&M\\
\midrule
\textbf{CRITERIOS \newline DE ACEPTACIÓN}& Se genera el documento de solicitud de permuta.\\
	\end{coolTable}
	\caption{Requisito de información 007\label{tab:RI-007}}
\end{table*}

\begin{table*}[htb]
	\centering
	\begin{coolTable}{p{4cm}p{\textwidth-6cm}}{2}{RI-008}
\textbf{DEFINICIÓN}& \textbf{Como} estudiante, \textbf{quiero} firmar la solicitud de permuta, {para} presentar el documento solicitado a la secretaría.\\
\midrule
\textbf{PRIORIDAD}&M\\
\midrule
\textbf{CRITERIOS \newline DE ACEPTACIÓN}& Se genera el documento firmado para la presentación en secretaría.\\
	\end{coolTable}
	\caption{Requisito de información 008\label{tab:RI-008}} 
\end{table*}

\begin{table*}[htb]
	\centering
	\begin{coolTable}{p{4cm}p{\textwidth-6cm}}{2}{RI-009}
\textbf{DEFINICIÓN}& {Como} estudiante, {quiero} consultar el estado de mis solicitudes de la gestión de permutas, \textbf {para} conocer si tengo que realizar alguna acción.\\
\midrule
\textbf{PRIORIDAD}&M\\
\midrule
\textbf{CRITERIOS \newline DE ACEPTACIÓN}& Se muestran todas las permutas que he realizado.\\
	\end{coolTable}
	\caption{Requisito de información 009\label{tab:RI-009}}
\end{table*}

\begin{table*}[htb]
	\centering
	\begin{coolTable}{p{4cm}p{\textwidth-6cm}}{2}{RI-010}
\textbf{DEFINICIÓN}& \textbf{Como} estudiante, \textbf{quiero} ver las notificaciones que se me han enviado, \textbf{para} poder continuar con el proceso que sea requerido.\\
\midrule
\textbf{PRIORIDAD}&M\\
\midrule
\textbf{CRITERIOS \newline DE ACEPTACIÓN}& Se muestran todas las notificaciones de un estudiante.\\
	\end{coolTable}
	\caption{Requisito de información 010\label{tab:RI-010}}
\end{table*}

\begin{table*}[htb]
	\centering
	\begin{coolTable}{p{4cm}p{\textwidth-6cm}}{2}{RI-011}
\textbf{DEFINICIÓN}& \textbf{Como} administrador, \textbf{quiero} visualizar las estadísticas de las permutas solicitadas, aceptadas, rechazadas \textbf{para} disponer de un panel en el que consultar de forma ágil datos relevantes del proceso.\\
\midrule
\textbf{PRIORIDAD}&M\\
\midrule
\textbf{CRITERIOS \newline DE ACEPTACIÓN}& Se visualizan las estadísticas referentes a los procesos de permuta registrados en el sistema.\\
	\end{coolTable}
	\caption{Requisito de información 011\label{tab:RI-011}}
\end{table*}

\begin{table*}[htb]
	\centering
	\begin{coolTable}{p{4cm}p{\textwidth-6cm}}{2}{RI-012}
\textbf{DEFINICIÓN}& \textbf{Como} administrador, \textbf{quiero} consultar los grados,cursos,grupos y asignaturas \textbf{para} verificar que toda la información disponible en el sistema es correcta.\\
\midrule
\textbf{PRIORIDAD}&M\\
\midrule
\textbf{CRITERIOS \newline DE ACEPTACIÓN}& Se visualizan los grados,cursos,grupos y asignaturas que se imparten en la escuela.\\
	\end{coolTable}
	\caption{Requisito de información 012\label{tab:RI-012}}
\end{table*}

\begin{table*}[htb]
	\centering
	\begin{coolTable}{p{4cm}p{\textwidth-6cm}}{2}{RI-013}
\textbf{DEFINICIÓN}& \textbf{Como} administrador, \textbf{quiero} gestionar las incidencias que puedan surgir durante el proceso \textbf{para} solucionar posibles errores y agilizar el proceso.\\
\midrule
\textbf{PRIORIDAD}&M\\
\midrule
\textbf{CRITERIOS \newline DE ACEPTACIÓN}& Se muestran las incidencias generadas por los estudiantes.\\
	\end{coolTable}
	\caption{Requisito de información 013\label{tab:RI-013}}
\end{table*}

\begin{table*}[htb]
	\centering
	\begin{coolTable}{p{4cm}p{\textwidth-6cm}}{2}{RI-014}
\textbf{DEFINICIÓN}& \textbf{Como} administrador, \textbf{quiero} ver las notificaciones que se me han enviado, \textbf{para} poder solucionar posibles errores.\\
\midrule
\textbf{PRIORIDAD}&M\\
\midrule
\textbf{CRITERIOS \newline DE ACEPTACIÓN}& Se muestran todas las notificaciones del sistema.\\
	\end{coolTable}
	\caption{Requisito de información 014\label{tab:RI-014}}
\end{table*}

\begin{table*}[htb]
	\centering
	\begin{coolTable}{p{4cm}p{\textwidth-6cm}}{2}{RI-015}
\textbf{DEFINICIÓN}& \textbf{Como} secretaría, \textbf{quiero} consultar las permutas que ha sido firmadas por ambas partes, \textbf{para} aceptar o rechazar según sea conveniente.\\
\midrule
\textbf{PRIORIDAD}&M\\
\midrule
\textbf{CRITERIOS \newline DE ACEPTACIÓN}& Se muestran todas las permutas que requieren que sean aceptadas o rechazadas.\\
	\end{coolTable}
	\caption{Requisito de información 015\label{tab:RI-015}}
\end{table*}


\begin{table*}[htb]
	\centering
	\begin{coolTable}{p{4cm}p{\textwidth-6cm}}{2}{RI-016}
\textbf{DEFINICIÓN}& \textbf{Como} secretaría, \textbf{quiero} modificar el estado a pendiente de correción, \textbf{para} que los datos sean correctos con respecto a los datos que tiene secretaría.\\
\midrule
\textbf{PRIORIDAD}&M\\
\midrule
\textbf{CRITERIOS \newline DE ACEPTACIÓN}& La posibilidad de modificar el estado de una permuta para que continue con el proceso.\\
	\end{coolTable}
	\caption{Requisito de información 016\label{tab:RI-016}}
\end{table*}

\begin{table*}[htb]
	\centering
	\begin{coolTable}{p{4cm}p{\textwidth-6cm}}{2}{RI-017}
\textbf{DEFINICIÓN}& \textbf{Como} secretaría, \textbf{quiero} poder aceptar o rechazar la permuta solicitada, \textbf{para} cumplir con la gestión requerida.\\
\midrule
\textbf{PRIORIDAD}&M\\
\midrule
\textbf{CRITERIOS \newline DE ACEPTACIÓN}& Poder conceder la permuta para concluir con el proceso.\\
	\end{coolTable}
	\caption{Requisito de información 017\label{tab:RI-017}}
\end{table*}

\begin{table*}[htb]
	\centering
	\begin{coolTable}{p{4cm}p{\textwidth-6cm}}{2}{RI-018}
\textbf{DEFINICIÓN}& \textbf{Como} secretaría, \textbf{quiero} ver las notificaciones que se han generado, \textbf{para} continuar con el proceso de la gestión permutas.\\
\midrule
\textbf{PRIORIDAD}&M\\
\midrule
\textbf{CRITERIOS \newline DE ACEPTACIÓN}& Se muestran las notificaciones generadas en el proceso de gestión de permutas.\\
	\end{coolTable}
	\caption{Requisito de información 018\label{tab:RI-018}}
\end{table*}



\clearpage
\subsection{Requisitos de calidad}

\begin{table*}[htb]
	\centering
	\begin{coolTable}{p{4cm}p{\textwidth-6cm}}{2}{RC-001}
\textbf{DEFINICIÓN}&\textbf{Como} usuario, \textbf{quiero} ver el las operaciones en la aplicación no tarden más de 10 segundos, \textbf {para} poder gestionar las permutas de manera veloz.\\
\midrule
\textbf{PRIORIDAD}&M\\
\midrule
\textbf{CRITERIOS \newline DE ACEPTACIÓN}& Todas las interacciones con la aplicación tienen un tiempo de respuesta menor o igual a 10 segundos.\\
	\end{coolTable}
	\caption{Requisito de calidad 001\label{tab:RC-001}}
\end{table*}

\clearpage
%-------------------------------------------------------------------------------------------------------
\subsection{Requisitos funcionales}

\begin{table*}[htb]
	\centering
	\begin{coolTable}{p{4cm}p{\textwidth-6cm}}{2}{RF-001}
\textbf{DEFINICIÓN}& \textbf {Como} usuario, \textbf {quiero} que las notificaciones siempre se encuentren visibles, \textbf {para} poder visualizar las notificaciones existentes.\\
\midrule
\textbf{PRIORIDAD}&M\\
\midrule
\textbf{CRITERIOS \newline DE ACEPTACIÓN}& Las notificaciones se encuentran visible en cada vista para consultarla cuando el usuario estime oportuno. \\
	\end{coolTable}
	\caption{Requisito funcional 001\label{tab:RF-001}}
\end{table*}

%-------------------------------------------------------------------------------------------------------

\clearpage

\subsection{Requisitos no funcionales}

\begin{table*}[htb]
	\centering
	\begin{coolTable}{p{4cm}p{\textwidth-6cm}}{2}{RNF-001}
\textbf{DEFINICIÓN}& \textbf {Como} administrador, \textbf {quiero} que el usuario tenga acceso a los términos del servicio (términos de uso), \textbf {para} que tengan constancia que significa utilizar nuestro servicio.\\
\midrule
\textbf{PRIORIDAD}&M\\
\midrule
\textbf{CRITERIOS \newline DE ACEPTACIÓN}& A la hora de iniciar sesión en la aplicación se muestran los términos de servicio. \\
	\end{coolTable}
	\caption{Requisito no funcional 001\label{tab:RNF-001}}
\end{table*}

%-------------------------------------------------------------------------------------------------------

\begin{table*}[htb]
	\centering
	\begin{coolTable}{p{4cm}p{\textwidth-6cm}}{2}{RNF-002}
\textbf{DEFINICIÓN}& \textbf {Como} administrador, \textbf {quiero} que el usuario tenga acceso al aviso de privacidad, \textbf {para} que el usuario sepa que tratamiento le daremos a sus datos.\\
\midrule
\textbf{PRIORIDAD}&M\\
\midrule
\textbf{CRITERIOS \newline DE ACEPTACIÓN}& A la hora de iniciar sesión en la aplicación se muestra la política de datos. \\
	\end{coolTable}
	\caption{Requisito no funcional 002\label{tab:RNF-002}}
\end{table*}

%-------------------------------------------------------------------------------------------------------

\begin{table*}[htb]
	\centering
	\begin{coolTable}{p{4cm}p{\textwidth-6cm}}{2}{RNF-003}
\textbf{DEFINICIÓN}& \textbf {Como} administrador, \textbf {quiero} que el usuario pueda tener acceso a los datos de la empresa, \textbf {para} que los usuarios tengan acceso a nuestros datos.\\
\midrule
\textbf{PRIORIDAD}&M\\
\midrule
\textbf{CRITERIOS \newline DE ACEPTACIÓN}& Podemos consultar los datos de la empresa en una pestaña de la aplicación. \\
	\end{coolTable}
	\caption{Requisito no funcional 003\label{tab:RNF-003}}
\end{table*}

%-------------------------------------------------------------------------------------------------------

\begin{table*}[htb]
	\centering
	\begin{coolTable}{p{4cm}p{\textwidth-6cm}}{2}{RNF-004}
\textbf{DEFINICIÓN}& \textbf {Como} administrador, \textbf {quiero} que el sistema funcione sobre un servidor web que opere con el protocolo HTTPS, \textbf {para} garantizar la seguridad de los usuarios.\\
\midrule
\textbf{PRIORIDAD}&M\\
\midrule
\textbf{CRITERIOS \newline DE ACEPTACIÓN}& Podemos comprobar si se cumple mirando la dirección del sitio web. \\
	\end{coolTable}
	\caption{Requisito no funcional 004\label{tab:RNF-004}}
\end{table*}

%-------------------------------------------------------------------------------------------------------